\documentclass[12pt,a4paper]{article}
\usepackage[utf8]{inputenc}
\usepackage[T1]{fontenc}
\usepackage[margin=1in]{geometry}
\usepackage{hyperref}
\usepackage{longtable}
\usepackage{booktabs}
\usepackage{array}
\usepackage{enumitem}
\usepackage{setspace}
\usepackage{xcolor}

\hypersetup{
  colorlinks=true,
  linkcolor=blue,
  urlcolor=blue,
  citecolor=blue,
  filecolor=blue
}

\onehalfspacing
\setlength{\parskip}{0.5em}
\setlength{\parindent}{0pt}

\begin{document}

\begin{center}
{\LARGE \textbf{CUROVA Healthcare Management System}}\\[0.3em]
{\large Final Project Documentation}
\end{center}

This document is the final consolidated project documentation, prepared by following the capstone guideline flow (deliverable sequence without timeline dependence), and reconciling early D1/D2 plans with the final implemented system.

\section{Deliverable-Based Development Flow (Guideline-Aligned)}

The project is documented in the standard development progression:
\begin{enumerate}
  \item Problem definition and scope
  \item Requirements analysis (functional and non-functional)
  \item System architecture and design decisions
  \item Data model and API design
  \item Implementation (backend and frontend)
  \item Testing and validation
  \item Deployment and operational setup
  \item Limitations, lessons learned, and future work
\end{enumerate}

This follows the expected capstone structure while intentionally ignoring date constraints.

\section{Problem Statement and Scope}

Healthcare workflows for appointment booking, medical record access, prescriptions, and lab workflows are often fragmented and role-dependent. Curova addresses this with a role-based web system for:

\begin{itemize}
  \item Patients: profile, appointments, records, prescriptions, documents, medications, lab results
  \item Doctors: schedule, patients, record creation, lab test review
  \item Lab technicians: lab order handling and result updates
  \item Admin: user and appointment oversight
\end{itemize}

The system is web-first, with REST APIs powering a React frontend.

\section{D1 and D2 vs Final Implementation}

\subsection{D1 (Initial Requirements) Baseline}

From \texttt{project\_docs/d1.tex}, core goals included:
\begin{itemize}
  \item Multi-role auth and access control
  \item Patient booking and history
  \item Doctor schedule and treatment workflow
  \item Admin user oversight
  \item Document handling
  \item Medication tracking and reminders
\end{itemize}

\subsection{D2 (Initial Design) Baseline}

From \texttt{project\_docs/d2.tex}, planned design emphasized:
\begin{itemize}
  \item Three-tier architecture (frontend, backend, database)
  \item Django + DRF backend with PostgreSQL
  \item Role-segmented UI and endpoint responsibilities
  \item ER-driven data model design
\end{itemize}

\subsection{Implementation Changes}

Compared with early D1/D2, final implementation evolved as follows:
\begin{enumerate}
  \item Stronger domain modularization: backend split into users, appointments, medical, medications, documents, lab\_tests, notifications, messaging.
  \item Notifications implemented as a dedicated domain.
  \item Messaging deferred: URL module exists but has no active endpoints.
  \item Medication reminder view logic exists, but reminder CRUD routes are not currently mounted in URL config.
  \item Deployment strategy pivoted to Render (frontend, backend, PostgreSQL).
\end{enumerate}

\section{Final System Architecture}

Curova uses a three-tier architecture:
\begin{enumerate}
  \item Presentation tier: React + Vite frontend in \texttt{frontend/src}
  \item Application tier: Django + DRF backend in \texttt{backend}
  \item Data tier: PostgreSQL relational database
\end{enumerate}

\subsection{Request Flow Summary}

A typical request flow:
\begin{enumerate}
  \item Frontend route/action in \texttt{frontend/src/App.jsx}
  \item Axios request from service layer (for example \texttt{frontend/src/services/api.js})
  \item Django URL resolution in \texttt{backend/curova\_backend/urls.py}
  \item DRF view performs auth and permission checks
  \item Serializer validates input/output
  \item ORM reads/writes PostgreSQL
  \item JSON response returns to frontend
\end{enumerate}

\section{Technology Stack}

\subsection{Backend}
\begin{itemize}
  \item Python
  \item Django 6.x
  \item Django REST Framework
  \item PostgreSQL
  \item DRF token authentication
  \item django-cors-headers
  \item gunicorn
\end{itemize}

Dependencies: \texttt{backend/requirements.txt}

\subsection{Frontend}
\begin{itemize}
  \item React 19
  \item Vite 7
  \item React Router
  \item Axios
\end{itemize}

Dependencies: \texttt{frontend/package.json}

\section{Backend Module Documentation}

\subsection{Root Configuration}
\begin{itemize}
  \item Global URLs: \texttt{backend/curova\_backend/urls.py}
  \item Settings: \texttt{backend/curova\_backend/settings.py}
  \item Pagination helper: \texttt{backend/curova\_backend/pagination.py}
\end{itemize}

\subsection{users app}
Key files:
\begin{itemize}
  \item \texttt{backend/users/models.py}
  \item \texttt{backend/users/views.py}
  \item \texttt{backend/users/permissions.py}
  \item \texttt{backend/users/urls.py}
  \item \texttt{backend/users/admin\_urls.py}
\end{itemize}
Responsibilities: registration, login/logout, Google login, profile and password updates, role-aware behavior, admin user and appointment oversight APIs.

\subsection{appointments app}
\begin{itemize}
  \item \texttt{backend/appointments/models.py}
  \item \texttt{backend/appointments/serializers.py}
  \item \texttt{backend/appointments/views.py}
  \item \texttt{backend/appointments/urls.py}
\end{itemize}
Responsibilities: booking flow, slot validation, role-scoped listing/detail, status updates, booked-slot queries.

\subsection{medical app}
\begin{itemize}
  \item \texttt{backend/medical/models.py}
  \item \texttt{backend/medical/serializers.py}
  \item \texttt{backend/medical/views.py}
  \item \texttt{backend/medical/urls.py}
\end{itemize}
Responsibilities: medical record read/write, prescription linkage, doctor-centric record creation.

\subsection{medications app}
\begin{itemize}
  \item \texttt{backend/medications/models.py}
  \item \texttt{backend/medications/serializers.py}
  \item \texttt{backend/medications/views.py}
  \item \texttt{backend/medications/urls.py}
  \item \texttt{backend/medications/reminders.py}
\end{itemize}
Responsibilities: medication listing/creation, reminder model support and helper logic. Reminder view methods exist but are not currently exposed in URL config.

\subsection{documents app}
\begin{itemize}
  \item \texttt{backend/documents/models.py}
  \item \texttt{backend/documents/serializers.py}
  \item \texttt{backend/documents/views.py}
  \item \texttt{backend/documents/urls.py}
\end{itemize}
Responsibilities: medical document upload/retrieval with ownership and role checks.

\subsection{lab\_tests app}
\begin{itemize}
  \item \texttt{backend/lab\_tests/models.py}
  \item \texttt{backend/lab\_tests/serializers.py}
  \item \texttt{backend/lab\_tests/views.py}
  \item \texttt{backend/lab\_tests/urls.py}
\end{itemize}
Responsibilities: lab order lifecycle, result upload/update, doctor/lab role-specific transitions.

\subsection{notifications app}
\begin{itemize}
  \item \texttt{backend/notifications/models.py}
  \item \texttt{backend/notifications/serializers.py}
  \item \texttt{backend/notifications/views.py}
  \item \texttt{backend/notifications/urls.py}
  \item \texttt{backend/notifications/helpers.py}
\end{itemize}
Responsibilities: list/unread count, mark read, mark all read, delete, clear, and event notifications from other domains.

\subsection{messaging app (deferred)}
\begin{itemize}
  \item \texttt{backend/messaging/urls.py} currently contains empty urlpatterns.
\end{itemize}

\section{Data Model Overview}

Primary entities:
\begin{enumerate}
  \item User
  \item Patient
  \item Doctor
  \item Appointment
  \item MedicalRecord
  \item Prescription
  \item Medication
  \item MedicationReminder
  \item MedicalDocument
  \item LabTest
  \item LabTestResult
  \item Notification
\end{enumerate}

Model references:
\texttt{backend/users/models.py},
\texttt{backend/appointments/models.py},
\texttt{backend/medical/models.py},
\texttt{backend/medications/models.py},
\texttt{backend/documents/models.py},
\texttt{backend/lab\_tests/models.py},
\texttt{backend/notifications/models.py}.

\section{Authentication and Authorization}

Authentication:
\begin{itemize}
  \item Token-based authentication using DRF token flow
  \item Primary auth endpoints in \texttt{backend/users/urls.py}
\end{itemize}

Authorization:
\begin{itemize}
  \item Role-aware permission classes in \texttt{backend/users/permissions.py}
  \item Endpoint-level checks inside app views
\end{itemize}

User roles:
\begin{itemize}
  \item patient
  \item doctor
  \item admin
  \item lab\_tech
\end{itemize}

\section{API Documentation (Backend)}

Base URL pattern:
\begin{itemize}
  \item Health: \texttt{/health/}
  \item API root: \texttt{/api/}
  \item Auth namespace: \texttt{/api/auth/}
  \item Admin namespace: \texttt{/api/admin/}
\end{itemize}

\subsection{Health Endpoint}

\begin{longtable}{p{2cm} p{3.6cm} p{6cm} p{2cm}}
\toprule
\textbf{Method} & \textbf{Path} & \textbf{Purpose} & \textbf{Auth} \\
\midrule
GET & /health/ & Service health check & No \\
\bottomrule
\end{longtable}

Defined in \texttt{backend/curova\_backend/urls.py}.

\subsection{Auth and User Endpoints}

\begin{longtable}{p{2cm} p{4.4cm} p{4.8cm} p{2.4cm}}
\toprule
\textbf{Method} & \textbf{Path} & \textbf{Purpose} & \textbf{Auth} \\
\midrule
POST & /api/auth/register/ & Register user & No \\
POST & /api/auth/login/ & Login and token issue & No \\
POST & /api/auth/google-login/ & Google login & No \\
POST & /api/auth/logout/ & Logout/token invalidation & Yes \\
POST & /api/auth/delete-account/ & Account anonymization/deactivation & Yes \\
GET, PUT & /api/auth/me/ & Current user profile details/update & Yes \\
POST & /api/auth/change-password/ & Password change & Yes \\
GET, PUT & /api/auth/profile/ & Patient profile view/update & Patient \\
GET & /api/auth/dashboard-stats/ & Patient dashboard stats & Patient \\
GET & /api/auth/doctors/ & Doctor list & Authenticated \\
GET & /api/auth/doctor/dashboard-stats/ & Doctor dashboard stats & Doctor \\
GET, PUT & /api/auth/doctor/profile/ & Doctor profile & Doctor \\
GET, PUT & /api/auth/doctor/schedule/ & Doctor schedule & Doctor \\
GET & /api/auth/doctor/patients/ & Doctor patient list & Doctor \\
GET & /api/auth/doctor/patients/{patient\_id}/ & Doctor patient detail & Doctor \\
\bottomrule
\end{longtable}

\subsection{Admin Endpoints}

\begin{longtable}{p{2cm} p{4.4cm} p{4.8cm} p{2.4cm}}
\toprule
\textbf{Method} & \textbf{Path} & \textbf{Purpose} & \textbf{Auth} \\
\midrule
GET & /api/admin/stats/ & Admin statistics & Admin \\
GET, POST & /api/admin/users/ & List/create users & Admin \\
GET, PATCH & /api/admin/users/{user\_id}/ & User detail/update state & Admin \\
GET & /api/admin/appointments/ & Appointment oversight list & Admin \\
GET & /api/admin/appointments/export/ & Appointment export & Admin \\
\bottomrule
\end{longtable}

\subsection{Appointments Endpoints}

\begin{longtable}{p{2cm} p{4.4cm} p{4.8cm} p{2.4cm}}
\toprule
\textbf{Method} & \textbf{Path} & \textbf{Purpose} & \textbf{Auth} \\
\midrule
GET, POST & /api/appointments/ & List appointments / create appointment & Authenticated (role-scoped) \\
GET & /api/appointments/booked-slots/ & Get booked slots by doctor/date & Authenticated \\
GET, PATCH & /api/appointments/{id}/ & Appointment detail / status update & Authenticated (ownership/role constrained) \\
\bottomrule
\end{longtable}

\subsection{Medical Record Endpoints}

\begin{longtable}{p{2cm} p{4.4cm} p{4.8cm} p{2.4cm}}
\toprule
\textbf{Method} & \textbf{Path} & \textbf{Purpose} & \textbf{Auth} \\
\midrule
GET & /api/records/ & List accessible medical records & Authenticated \\
GET & /api/records/{id}/ & Medical record detail & Authenticated \\
POST & /api/records/create/ & Create medical record & Doctor \\
\bottomrule
\end{longtable}

\subsection{Medication Endpoints}

\begin{longtable}{p{2cm} p{4.4cm} p{4.8cm} p{2.4cm}}
\toprule
\textbf{Method} & \textbf{Path} & \textbf{Purpose} & \textbf{Auth} \\
\midrule
GET & /api/medications/ & List active medications & Patient \\
\bottomrule
\end{longtable}

Note: Reminder endpoints are implemented in views but not exposed in current URL config.

\subsection{Document Endpoints}

\begin{longtable}{p{2cm} p{4.4cm} p{4.8cm} p{2.4cm}}
\toprule
\textbf{Method} & \textbf{Path} & \textbf{Purpose} & \textbf{Auth} \\
\midrule
GET, POST & /api/documents/ & List documents / upload document & Authenticated (role checks) \\
GET, DELETE & /api/documents/{id}/ & Document detail / delete & Authenticated (ownership/role checks) \\
\bottomrule
\end{longtable}

\subsection{Lab Test Endpoints}

\begin{longtable}{p{2cm} p{4.4cm} p{4.8cm} p{2.4cm}}
\toprule
\textbf{Method} & \textbf{Path} & \textbf{Purpose} & \textbf{Auth} \\
\midrule
GET, POST & /api/lab-tests/ & List/order lab tests & Authenticated (role constrained) \\
GET, PATCH & /api/lab-tests/{id}/ & Lab test detail / status update & Authenticated (role constrained) \\
GET, POST & /api/lab-tests/{id}/result/ & View or create test result & Authenticated (role constrained) \\
PATCH & /api/lab-tests/{id}/result/update/ & Update existing test result & Authenticated (role constrained) \\
\bottomrule
\end{longtable}

\subsection{Notification Endpoints}

\begin{longtable}{p{2cm} p{4.4cm} p{4.8cm} p{2.4cm}}
\toprule
\textbf{Method} & \textbf{Path} & \textbf{Purpose} & \textbf{Auth} \\
\midrule
GET & /api/notifications/ & List notifications & Authenticated \\
PATCH & /api/notifications/{id}/read/ & Mark one notification read & Authenticated \\
DELETE & /api/notifications/{id}/delete/ & Delete one notification & Authenticated \\
POST & /api/notifications/mark-all-read/ & Mark all notifications read & Authenticated \\
GET & /api/notifications/unread-count/ & Get unread notification count & Authenticated \\
DELETE & /api/notifications/clear/ & Clear all notifications & Authenticated \\
\bottomrule
\end{longtable}

\section{Frontend Documentation Summary}

Routing and role-segmented page architecture are defined in \texttt{frontend/src/App.jsx}.

Route groups:
\begin{itemize}
  \item Public: homepage, login, registration, policy pages
  \item Patient: dashboard, appointments, records, prescriptions, documents, lab results, medications, notifications, profile
  \item Doctor: dashboard, schedule, patients, patient detail, create record, notifications, profile
  \item Admin: dashboard, users, appointments, notifications
  \item Lab tech: orders, notifications, profile
\end{itemize}

\section{Testing and Validation Status}

Current status is primarily manual validation and integration checks.

Evidence in repository:
\begin{itemize}
  \item Backend test placeholders exist in each app test file under \texttt{backend}
  \item Functional checks were performed through manual role-based flows and deployment checks
\end{itemize}

Improvement needed:
\begin{enumerate}
  \item Unit tests for serializers and validators
  \item API integration tests by role
  \item Permission boundary tests
  \item End-to-end smoke tests for key user journeys
\end{enumerate}

\section{Deployment Documentation (Render)}

Current deployment model:
\begin{itemize}
  \item Frontend on Render static service
  \item Backend on Render web service (Gunicorn)
  \item PostgreSQL on Render managed database
\end{itemize}

Operational references:
\begin{itemize}
  \item Root Render blueprint: \texttt{render.yaml}
  \item Backend environment-based settings: \texttt{backend/curova\_backend/settings.py}
\end{itemize}

Production essentials:
\begin{enumerate}
  \item Correct environment variables (ALLOWED\_HOSTS, CORS\_ALLOWED\_ORIGINS, DATABASE\_URL, secrets)
  \item Migrations on deploy
  \item Health endpoint monitoring at \texttt{/health/}
\end{enumerate}

\section{Known Limitations and Risks}

\begin{enumerate}
  \item Messaging feature is not exposed yet (\texttt{backend/messaging/urls.py}).
  \item Medication reminder endpoint exposure gap: reminder view logic exists but routing is incomplete in \texttt{backend/medications/urls.py}.
  \item Automated test coverage is limited.
  \item Advanced UX/ops hardening remains pending.
\end{enumerate}

\section{Future Work Roadmap}

\subsection*{Short Term}
\begin{enumerate}
  \item Expose reminder CRUD routes and complete reminder UI workflow
  \item Complete messaging API and frontend feature
  \item Add comprehensive backend automated tests
  \item Improve audit and reporting capabilities in admin
\end{enumerate}

\subsection*{Medium Term}
\begin{enumerate}
  \item Event-driven notification pipeline improvements
  \item Enhanced observability (metrics, dashboards, alerts)
  \item Scalability and performance profiling for high-load flows
\end{enumerate}

\subsection*{Long Term}
\begin{enumerate}
  \item Deeper compliance hardening and audit-trail extension
  \item Mobile client integration over existing API surface
\end{enumerate}

\section{Appendix: Key Project References}

Core design and planning references:
\begin{itemize}
  \item Guideline source: \texttt{project\_docs/capstone\_project\_guidelines.pdf}
  \item D1 baseline: \texttt{project\_docs/D1\_requirement\_doc.pdf}, \texttt{project\_docs/d1.tex}
  \item D2 baseline: \texttt{project\_docs/D2 Design Decisions (1).pdf}, \texttt{project\_docs/d2.tex}
  \item Planning decisions: \texttt{project\_docs/PLAN.md}, \texttt{project\_docs/FINAL\_DECISIONS.md}
\end{itemize}

Interview-oriented references in repository:
\begin{itemize}
  \item \texttt{BACKEND\_INTERVIEW\_GUIDE.md}
  \item \texttt{BACKEND\_FOUNDATIONS\_GUIDE.md}
\end{itemize}

\end{document}
